%
%  chapter2.tex  —  Formulação do Problema e Implementação
%
\chapter{Formulação do Problema e Implementação}
\label{cha:formulacao}

\section{Representação do Estado}
\label{sec:representacao}

\subsection{Sudoku Clássico}
\label{ssec:rep_classico}

A grelha $9 \times 9$ é representada como uma lista plana de 81 inteiros, indexada linha a linha (da esquerda para a direita, de cima para baixo). O índice plano de uma célula na linha $r$ e coluna $c$ (ambas de base 1) é dado por $(r-1) \times 9 + (c-1)$. Uma célula vazia é representada pelo valor $0$; as células preenchidas pelo puzzle (\emph{givens}) têm valores de 1 a 9.

Esta representação plana permite calcular as coordenadas de linha, coluna e caixa sem estruturas de dados auxiliares. O índice da caixa de uma célula $(r, c)$ é $\lfloor(r-1)/3\rfloor \times 3 + \lfloor(c-1)/3\rfloor$.

Em Prolog, os predicados fundamentais são:
\begin{itemize}
    \item \texttt{all\_different/1} — verifica que uma lista não contém duplicados;
    \item \texttt{valid/1} — verifica que todas as linhas, colunas e caixas satisfazem \texttt{all\_different};
    \item \texttt{candidates/3} — dada uma célula vazia, calcula os valores 1--9 que não violam nenhuma restrição da célula;
    \item \texttt{next\_empty/2} — selecciona a célula vazia com menos candidatos (MRV).
\end{itemize}

Em Python (\texttt{board.py}), a representação é equivalente, com operações como \texttt{load\_board/1} para ler o puzzle de ficheiro e \texttt{format\_board/1} para imprimir a grelha.

\subsection{Killer Sudoku}
\label{ssec:rep_killer}

O Killer Sudoku partilha a mesma representação de grelha. A informação adicional das gaiolas é guardada como uma lista de pares $(s, C)$, onde $s$ é a soma alvo e $C$ é a lista de células (como pares linha-coluna de base 1) que formam a gaiola.

Em Prolog, as gaiolas são codificadas como factos \texttt{killer\_cage(PuzzleId, Sum, Cells)} no ficheiro \texttt{killer\_sudoku.pl}. Em Python, o formato de ficheiro é:
\begin{verbatim}
  cage(Sum, [(R1,C1),(R2,C2),...]).
\end{verbatim}
que é interpretado pela função \texttt{load\_killer\_puzzle} em \texttt{board.py}.

\section{Pesquisa em Espaço de Estados — Prolog}
\label{sec:prolog}

Os algoritmos de pesquisa em Prolog operam sobre \emph{atribuições parciais}: começam pelo puzzle inicial (com os \emph{givens} fixos e as restantes células a zero) e atribuem progressivamente valores às células vazias. Cada nó no espaço de estados é uma grelha parcialmente preenchida que satisfaz as restrições até ao momento.

\subsection{Aprofundamento Iterativo (DFID)}
\label{ssec:dfid}

O aprofundamento iterativo com pesquisa em profundidade (DFID) é um algoritmo de pesquisa não informada que combina a completude e optimalidade da pesquisa em largura com a eficiência em memória da pesquisa em profundidade \cite{russell2020ai}. O algoritmo executa sucessivas iterações de pesquisa em profundidade limitada, aumentando o limite de profundidade em cada iteração.

\paragraph{Limite de Profundidade com MRV.}
No contexto do Sudoku, o limite de profundidade natural é o número de células vazias no puzzle: cada nó interno corresponde a uma célula preenchida, pelo que a solução está exactamente a essa profundidade. Quando se usa MRV para seleccionar a próxima célula a preencher, a pesquisa é guiada pelas células mais restringidas primeiro, o que tipicamente reduz o ramo para um único filho válido (ou um número muito pequeno). Como consequência, a iteração correspondente ao limite máximo (número de células vazias) é a única que atinge nós folha, e as iterações anteriores terminam rapidamente por não encontrar folhas. O DFID \emph{degenera efectivamente numa DFS com constraint checking}, executando uma única passagem significativa.

\paragraph{Implementação.}
O predicado principal é \texttt{solve\_file/2}, que lê o puzzle de um ficheiro, inicializa um contador de nós e chama \texttt{dfid/4}:
\begin{verbatim}
  dfid(Board, Depth, MaxDepth, Nodes) :-
      Depth =< MaxDepth,
      ( solved(Board) -> true
      ; next_empty(Board, Pos),
        candidates(Board, Pos, Vals),
        member(Val, Vals),
        place(Board, Pos, Val, Board2),
        Nodes1 is Nodes + 1,
        dfid(Board2, Depth+1, MaxDepth, Nodes1)
      ).
\end{verbatim}

\subsection{Pesquisa A*}
\label{ssec:astar}

A pesquisa A* é um algoritmo de pesquisa informada que expande o nó com menor valor de $f(n) = g(n) + h(n)$, onde $g(n)$ é o custo do caminho desde o estado inicial até $n$ e $h(n)$ é uma estimativa heurística do custo de $n$ até ao objectivo \cite{russell2020ai}.

\paragraph{Heurística Admissível.}
A heurística adoptada é $h(n) = $ número de células vazias no estado $n$. Esta heurística é \emph{admissível} porque, para resolver cada célula vazia, é necessário efectuar pelo menos uma atribuição, pelo que $h(n)$ nunca sobrestima o custo real. Formalmente, se a solução se encontra a $k$ células de profundidade a partir de $n$, então $h(n) = k = $ custo real, pelo que $h(n) \leq h^*(n)$.

\paragraph{Comportamento com esta Heurística.}
Com $g(n) = $ número de células já preenchidas e $h(n) = $ número de células vazias, tem-se $f(n) = g(n) + h(n) = 81 - (\text{givens})$, que é \emph{constante} para todos os nós ao mesmo nível de profundidade. Neste caso, o A* comporta-se essencialmente como uma pesquisa em largura ao nível da ordenação da fila de prioridade, e a poda advém inteiramente da verificação de restrições (candidatos vazios), não da ordenação heurística.

\paragraph{Implementação.}
O A* usa uma fila de prioridade implementada com a biblioteca \texttt{library(heaps)} do SWI-Prolog. Cada elemento da fila é um par $(f, \text{Board})$. Uma lista fechada mantém os estados já expandidos. A expansão de cada nó usa \texttt{next\_empty/2} (MRV) para seleccionar a célula a preencher:
\begin{verbatim}
  astar(OpenHeap, Closed, Solution, Nodes) :-
      get_from_heap(OpenHeap, _, Board, Rest),
      ( solved(Board) -> Solution = Board
      ; \+ member(Board, Closed),
        next_empty(Board, Pos),
        candidates(Board, Pos, Vals),
        expand_children(Board, Pos, Vals, Rest,
                        [Board|Closed], NewOpen, Nodes1),
        astar(NewOpen, [Board|Closed], Solution, Nodes1)
      ).
\end{verbatim}

\subsection{Extensão para Killer Sudoku — Prolog}
\label{ssec:killer_prolog}

A extensão para Killer Sudoku em Prolog (\texttt{killer\_sudoku.pl}) acrescenta dois predicados de verificação de gaiolas:

\begin{itemize}
    \item \texttt{valid\_cage(Board, cage(Sum, Cells))} — verifica que, quando todas as células da gaiola estão preenchidas, a sua soma é igual a \texttt{Sum} e não existem repetições;
    \item \texttt{partial\_cage\_ok(Board, cage(Sum, Cells))} — verifica a consistência parcial durante a pesquisa: se já existem valores parcialmente atribuídos na gaiola, a sua soma parcial não pode exceder \texttt{Sum}, e os valores já atribuídos não se podem repetir.
\end{itemize}

O predicado \texttt{partial\_cage\_ok} é crucial para a poda: ao ser chamado após cada atribuição durante a DFS/A*, elimina ramos que já violam as restrições das gaiolas, mesmo antes de a gaiola estar completamente preenchida.

Os solucionadores \texttt{solve\_dfid\_killer/3} e \texttt{solve\_astar\_killer/3} são extensões directas dos solucionadores clássicos, com a verificação adicional de \texttt{partial\_cage\_ok} em cada passo de expansão.

\section{Pesquisa em Espaço de Soluções — Python}
\label{sec:python}

Ao contrário dos algoritmos de pesquisa em espaço de estados, o SA e o GA operam sobre \emph{soluções completas}: começam com uma grelha $9 \times 9$ completamente preenchida (possivelmente inválida) e aplicam operadores que a transformam noutra solução completa, tentando minimizar o número de violações de restrições.

\subsection{Inicialização e Operador}
\label{ssec:init_op}

\paragraph{Inicialização.}
Cada caixa $3 \times 3$ é preenchida independentemente com os valores em falta (1--9 menos os \emph{givens} da caixa), dispostos de forma aleatória. Esta inicialização garante que \emph{as restrições de caixa estão sempre satisfeitas por construção}. Apenas as restrições de linha e coluna (e de gaiola, no Killer) podem ser violadas.

\paragraph{Operador de Troca Intra-caixa.}
O operador de vizinhança selecciona uma caixa aleatória e troca dois valores mutáveis (não-\emph{given}) dentro dela. Como a troca ocorre dentro da mesma caixa, as restrições de caixa continuam satisfeitas. Formalmente, se a caixa contém os valores $\{v_1, \ldots, v_9\}$ antes da troca, continua a contê-los depois.

\paragraph{Função de Custo.}
O custo $C(\text{board})$ conta as violações de linha e coluna:
\begin{equation}
    C = \sum_{i=1}^{9} (9 - |\text{distinct}(\text{row}_i)|) + \sum_{j=1}^{9} (9 - |\text{distinct}(\text{col}_j)|)
\end{equation}
O custo mínimo é 0 (solução completa e válida) e o máximo teórico é 144. No Killer Sudoku, soma-se adicionalmente o custo de gaiola:
\begin{equation}
    C_{\text{cage}} = \sum_{\text{gaiola}} \mathbb{1}[\text{soma} \neq \text{alvo}] + (\text{tamanho} - |\text{distinct}(\text{gaiola})|)
\end{equation}

\subsection{Arrefecimento Simulado (SA)}
\label{ssec:sa}

O arrefecimento simulado (SA) \cite{kirkpatrick1983optimization} é uma metaheurística que aceita soluções piores com uma probabilidade que diminui ao longo do tempo, permitindo escapar de mínimos locais.

\paragraph{Esquema de Temperatura.}
A temperatura decresce de forma geométrica: $T_{t+1} = T_t \times \alpha$, onde $\alpha = 0.9999$ é o factor de arrefecimento. Com $T_{\text{init}} = 2.0$ e $T_{\text{min}} = 0.001$, cada execução natural tem $\approx 76\,000$ iterações.

\paragraph{Aceitação de Soluções.}
Um estado vizinho com custo $C_{\text{viz}}$ é aceite se $C_{\text{viz}} < C_{\text{actual}}$ (melhoria) ou com probabilidade $e^{-\Delta C / T}$ se $\Delta C = C_{\text{viz}} - C_{\text{actual}} > 0$ (aceitação probabilística). Esta regra permite escapar de mínimos locais quando a temperatura é elevada.

\paragraph{Reinicialização por Conclusão Natural.}
Quando a temperatura atinge $T_{\text{min}}$ sem encontrar a solução, o tabuleiro é reinicializado aleatoriamente e a temperatura é reposta em $T_{\text{init}}$. Esta estratégia de \emph{reinicialização por conclusão natural} garante que cada execução percorre o ciclo completo de arrefecimento, da fase exploratória (alta temperatura) à fase exploitativa (baixa temperatura), antes de recomeçar. O número total de iterações é limitado a $300\,000$.

\paragraph{Mínimo Local a Custo 2.}
Em puzzles de elevada dificuldade, o SA pode ficar preso num estado com custo 2 — duas violações em linhas ou colunas diferentes que envolvem células em caixas distintas. O operador de troca intra-caixa não pode resolver este problema directamente: corrigir uma violação requer trocar células entre caixas diferentes, o que violaria as restrições de caixa. A baixa temperatura torna a probabilidade de aceitar estados com custo superior praticamente nula ($e^{-1/0.001} \approx 0$), pelo que o algoritmo fica bloqueado. A solução é aguardar pelo arrefecimento completo e reinicializar.

\subsection{Algoritmo Genético (GA)}
\label{ssec:ga}

O algoritmo genético (GA) \cite{holland1975adaptation} mantém uma população de soluções candidatas e aplica operadores inspirados na evolução biológica (selecção, cruzamento, mutação) para melhorar a população ao longo das gerações.

\paragraph{Cromossoma.}
Cada cromossoma é uma grelha $9 \times 9$ inicializada da mesma forma que o SA (restrições de caixa garantidas por construção).

\paragraph{Função de Fitness.}
$f(\text{cromossoma}) = -C(\text{cromossoma})$, onde $C$ é a função de custo definida na Secção~\ref{ssec:init_op}. O GA \emph{maximiza} o fitness, equivalente a minimizar o custo.

\paragraph{Selecção por Torneio.}
São seleccionados $k = 3$ indivíduos aleatórios da população e o de maior fitness é escolhido como progenitor. Este processo é repetido para seleccionar dois progenitores ($p_1, p_2$) para o cruzamento.

\paragraph{Cruzamento Preservador de Caixas.}
Para cada uma das 9 caixas, o filho herda a caixa inteira do $p_1$ ou do $p_2$ com probabilidade 0.5. As células \emph{given} permanecem fixas. Este operador garante que as restrições de caixa se mantêm no filho.

\paragraph{Mutação.}
Com probabilidade $p_m = 0.25$, aplica-se a um filho o operador de troca intra-caixa (igual ao do SA), realizando 2 trocas.

\paragraph{Elitismo e Reinicialização por Estagnação.}
Os 2 melhores indivíduos da geração actual sobrevivem inalterados para a próxima geração (\emph{elitismo}). Se o melhor fitness não melhorar durante $2000$ gerações consecutivas (\emph{estagnação}), a população é reinicializada aleatoriamente, mantendo apenas o melhor indivíduo actual.

\subsection{Extensão para Killer Sudoku — Python}
\label{ssec:killer_python}

A extensão para Killer Sudoku em Python é transparente para os algoritmos SA e GA: a função de custo é aumentada com o custo de gaiola $C_{\text{cage}}$ (Secção~\ref{ssec:init_op}), implementado em \texttt{board.py::cage\_cost}. O carregamento do puzzle é tratado por \texttt{load\_killer\_puzzle}, que lê a grelha e as definições de gaiolas do ficheiro de texto. Tanto o SA como o GA seleccionam automaticamente o modo Killer quando o puzzle fornecido pertence ao conjunto \texttt{KILLER\_PUZZLES}.
